% generated by roam.py handbook export -- не редактировать, править заметки
% 20261007161134-условия_армихо_вульфа.org (body 6822678aba64)

\begin{definition}[{Условия Армихо-Вульфа}]\label[definition]{note:5330B23D-95DD-4D5D-B4A3-63EBADB1DB7D}
  Пусть $\phi: \R \rightarrow \R$, $\phi \in C^1(\R)$~--- непрерывно \hyperref[note:957FA09D-8883-40A5-B4C7-237D9F815A42]{дифференцируема}, рассмотрим отрезок $[a;b]$, такой что $\phi^{\prime}(a) < 0$

  Зафиксируем $0<c_1 <1$, $0 < c_2 < 1$;
  Ищем шаг $\alpha_k > a$, для которого выполнены оба условия:

  \hyperref[note:15611E05-13FC-4FB0-AE00-A4BB857A29E6]{Условие Армихо}:
  \[
  \phi(\alpha_k) <  \phi(a) + c_1(\alpha_k - a) \cdot \phi^{\prime}(a)
  \]

  \hyperref[note:2084FD2F-E2AA-4081-B9D1-14CEA27EA243]{Условие Вульфа}:
  \[
  | \phi^{\prime}(\alpha_k) | < c_2|\phi^{\prime}(a)|
  \]

  Такой шаг существует при $c_1 < c_2$~--- \hyperref[note:C143295F-C422-4FEE-95AD-3D6B620BC0E8]{Существование шага для Армихо-Вульфа}.
\end{definition}

\begin{remark}
  По сути Условие Армихо ограничивает нас от того, чтобы уйти слишком далеко от локального минимума (на другую ветку). А условие Вульфа ограничивает нас от того, чтобы делать слишком маленькие шаги.
\end{remark}

\begin{remark}
  Условие с модулем называют \emph{сильным} условием Вульфа; слабое~--- $\phi^{\prime}(\alpha_k) \ge c_2 \phi^{\prime}(a)$ (запрещает только слишком крутой спуск в новой точке, но не подъём).
\end{remark}
